% =====================================================================
%  Capítulo 4 — Materiais e Métodos
% =====================================================================
\chapter{Materiais e Métodos}
\label{cap:metodologia}

Este capítulo descreve a natureza da pesquisa, o processo de engenharia
adotado, os materiais e o ambiente utilizados e os critérios de validação,
permitindo a reprodução ou a avaliação crítica do trabalho. Procura-se
distinguir claramente o que foi \emph{planejado} (nos artefatos de
especificação e design) do que foi \emph{executado} (código e resultados).

\section{Natureza e abordagem da pesquisa}

O trabalho caracteriza-se como \textbf{pesquisa aplicada} com componente de
\textbf{desenvolvimento experimental}: parte de um problema de engenharia real
(volume de dados e necessidade de sinais de defeito no monitoramento de
condição) e entrega um artefato funcional --- um simulador --- com validação
por testes. A abordagem é essencialmente tecnológica e combina três áreas:
engenharia de software, processamento digital de sinais e engenharia de
confiabilidade (manutenção preditiva).

\section{Processo de engenharia: especificação dirigida}

O desenvolvimento seguiu o processo de \emph{Spec-Driven Development} (SDD),
no qual o código é consequência de uma especificação rastreável, e não o
contrário. O fluxo executado foi:

\begin{enumerate}
  \item \textbf{Auditoria do contexto}: análise das diretrizes técnicas de
  origem (critérios metrológicos, critério \(R^3\), catálogo de defeitos),
  disponíveis como documentos de entrada do projeto.
  \item \textbf{Especificação}: formalização dos requisitos funcionais
  (FR-001 a FR-019), requisitos não funcionais (NFR-001 a NFR-007) e critérios
  de aceitação (CA-001 a CA-014) em \texttt{SPECIFICATION.md}, com matriz de
  rastreabilidade.
  \item \textbf{Design}: arquitetura, contratos de API, modelo de dados e
  detalhamento por tarefas pequenas (T-001 a T-013) em \texttt{.specs/}.
  \item \textbf{Implementação}: execução das tarefas com gate mínimo de
  lint, type-check e testes por módulo.
  \item \textbf{Validação}: testes unitários, de integração e de contrato, e
  ensaio de ponta a ponta no navegador.
\end{enumerate}

Cada etapa deixou evidência verificável (especificação, design, tarefas,
testes e registros de execução), o que é utilizado neste texto como base das
afirmações sobre o sistema.

\section{Materiais e ambiente}

O sistema foi desenvolvido e validado no ambiente descrito na
Tabela~\ref{tab:ferramentas}. As versões informadas correspondem às
declaradas nos artefatos do projeto e verificadas durante a execução.

\begin{table}[htbp]
\centering
\caption{Materiais, ferramentas e versões utilizados no desenvolvimento e na validação.}
\label{tab:ferramentas}
\begin{tabular}{@{}lll@{}}
\toprule
\textbf{Camada} & \textbf{Ferramenta} & \textbf{Versão/nota} \\
\midrule
Back-end & Python & 3.12 \\
         & FastAPI & $\ge$ 0.115 (Pydantic v2) \\
         & NumPy / SciPy & $\ge$ 1.26 / $\ge$ 1.13 \\
         & SQLAlchemy + asyncpg + Alembic & 2.x / $\ge$ 0.29 / $\ge$ 1.13 \\
         & ruff / mypy / pytest & lint, format, type-check, testes \\
Front-end & React & 19.2.x instalada \\
          & Vite / TypeScript & 8.2.x / 6.0.x \\
          & Recharts & 3.10.x \\
          & oxlint / Vitest / Testing Library & lint e testes \\
Banco & PostgreSQL & 16 (contêiner Docker) \\
Infra & Docker / Docker Compose & v2 \\
Execução & Navegador & ensaio E2E em localhost:5173 \\
\bottomrule
\end{tabular}
\par\vspace{2pt}{\footnotesize Fonte: \texttt{frontend/package.json}, \texttt{backend/requirements*.txt} e \texttt{docker-compose.yml} do projeto.}
\end{table}

\section{Arquitetura de teste em camadas}

A estratégia de validação seguiu uma pirâmide de testes, com camadas de custo
crescente e escopo decrescente:

\begin{enumerate}
  \item \textbf{Testes unitários} dos módulos puros de domínio
  (\texttt{signal\_generator}, \texttt{fft\_processor}, \texttt{discard\_engine}), sem acesso a banco
  de dados ou rede. O motor de descarte recebe a última leitura persistida como
  parâmetro injetado, o que permite cobrir todas as combinações da regra de ouro
  e do Paralelepípedo de Descarte sem infraestrutura.
  \item \textbf{Testes de integração} da camada de persistência, executados
  contra um PostgreSQL real provisório (contêiner via \texttt{testcontainers}),
  validando migração, transações e o ponteiro de última leitura persistida.
  \item \textbf{Testes de contrato} da API, usando cliente HTTP assíncrono
  contra a aplicação FastAPI, verificando os contratos de \texttt{POST /simulacoes},
  \texttt{POST /snapshots} e \texttt{GET /health}, incluindo validações de entrada e a
  violação de Nyquist.
  \item \textbf{Testes de front-end} (componentes e validação de formulário)
  com Vitest e Testing Library.
  \item \textbf{Ensaio de ponta a ponta} no navegador, contra o ambiente
  Docker Compose completo, verificando o fluxo real de simulação e a renderização
  dos elementos visuais.
\end{enumerate}

Os critérios de aceitação funcionais (CA-001 a CA-014) foram usados como
fonte dos casos de teste, ligando cada teste ao requisito que implementa.

\section{Métricas e critérios de validação}

As métricas observadas na validação foram:

\begin{itemize}
  \item \textbf{correção espectral}: presença e ordem das componentes
  esperadas para cada defeito do catálogo, verificadas no espectro retornado
  pela API;
  \item \textbf{correção do descarte}: decisão (persistir/descartar) de acordo
  com a regra de ouro e o Paralelepípedo \(R^3\), conforme os casos de teste;
  \item \textbf{tempo de processamento} por simulação (indicador exposto pela
  API, FR-017);
  \item \textbf{taxa de descarte acumulada} (leituras descartadas divididas
  pelas leituras avaliadas, NFR-004);
  \item \textbf{gates de qualidade}: lint e type-check limpos e suíte de
  testes verde nas duas camadas.
\end{itemize}

\section{Protocolo experimental executado}

O protocolo efetivamente executado foi o seguinte (distinguindo-se o planejado
do executado):

\begin{enumerate}
  \item subir o ambiente completo com \texttt{start.sh} (Docker Compose: banco,
  migrações Alembic, back-end e front-end);
  \item semear um ponto de medição de demonstração com identificador fixo;
  \item executar a suíte de testes do back-end (lint + type-check + pytest) e
  do front-end (type-check + lint + vitest);
  \item verificar, via chamadas reais à API, as assinaturas espectrais de
  defeitos representativos (ex.: desbalanceamento com \(1X\) dominante,
  desalinhamento paralelo com \(2X\)/\(4X\)/\(6X\), \emph{oil whirl} com
  razão subsíncrona e rolamento com BPFI e \emph{sidebands});
  \item executar o ensaio de ponta a ponta no navegador, simulando um defeito
  e verificando os elementos visuais do painel;
  \item registrar as evidências (contagens de testes, capturas de tela) para
  este texto.
\end{enumerate}

O item 2 (seed de ponto) foi executado por script de inicialização, uma vez que
o schema não expõe endpoint de criação de hierarquia no MVP; essa limitação é
discutida no Capítulo~\ref{cap:discussao}.

\section{Limitações metodológicas}

As principais limitações metodológicas são: (i) a validação é local e não há
ambiente de produção definido (NFR-005); (ii) a avaliação de desempenho
baseou-se em testes funcionais e observação direta, sem bateria estatística
formal de medições; e (iii) não foi conduzida revisão sistemática da literatura,
apenas revisão delimitada. Essas limitações estão registradas nas pendências do
trabalho e devem ser tratadas com o orientador.
